% Exercices type bac — Variables aléatoires — v28

\subsection*{Sujet 1 — Jeu de jetons et espérance}

Une urne contient 10 jetons indiscernables au toucher : 1 jeton rouge, 3 jetons bleus et 6 jetons blancs.
Un joueur tire un jeton au hasard. Le jeu coûte 2 euros.
S'il tire le jeton rouge, il reçoit 8 euros. S'il tire un jeton bleu, il reçoit 3 euros. S'il tire un jeton blanc, il ne reçoit rien.
On note $X$ le gain algébrique du joueur, c'est-à-dire la somme reçue moins le prix du jeu.

\begin{enumerate}
\item Déterminer les valeurs prises par $X$.
\item Déterminer la loi de probabilité de $X$ sous forme de tableau.
\item Calculer l'espérance de $X$.
\item Interpréter ce résultat dans le contexte.
\item Le jeu est-il favorable au joueur ?
\item Quel prix d'entrée faudrait-il choisir pour que le jeu soit équitable ?
\end{enumerate}

\subsection*{Sujet 2 — Deux jeux à comparer}

On propose deux jeux. Dans chaque jeu, on lance un dé équilibré à six faces.

Jeu A : si le résultat est 6, on gagne 9 euros ; si le résultat est 4 ou 5, on gagne 3 euros ; sinon, on perd 3 euros.

Jeu B : si le résultat est pair, on gagne 3 euros ; si le résultat est impair, on perd 3 euros.

On note $A$ la variable aléatoire donnant le gain au jeu A et $B$ la variable aléatoire donnant le gain au jeu B.

\begin{enumerate}
\item Déterminer la loi de probabilité de $A$.
\item Déterminer la loi de probabilité de $B$.
\item Calculer $E(A)$ et $E(B)$.
\item Calculer $V(A)$ et $V(B)$.
\item En déduire les écarts types de $A$ et $B$.
\item Les deux jeux ont-ils le même gain moyen ?
\item Quel jeu paraît le plus risqué ? Justifier.
\end{enumerate}

\subsection*{Sujet 3 — Remise commerciale et prix payé}

Un magasin organise une opération promotionnelle. À chaque passage en caisse, le client tire une carte au hasard parmi 20 cartes indiscernables.
Deux cartes donnent une remise de 20 euros, cinq cartes donnent une remise de 10 euros, et les autres cartes ne donnent aucune remise.
Un client achète un article coûtant 50 euros. On note $P$ la variable aléatoire égale au prix réellement payé par le client après remise.

\begin{enumerate}
\item Déterminer les valeurs prises par $P$.
\item Déterminer la loi de probabilité de $P$.
\item Calculer $E(P)$.
\item Interpréter $E(P)$ dans le contexte.
\item Calculer $V(P)$.
\item Calculer l'écart type de $P$.
\item Le magasin affirme : « en moyenne, chaque client bénéficie de 5 euros de remise ». Cette affirmation est-elle vraie ?
\end{enumerate}
